% =============================================================================
% SEZIONE: Contesto Storico Espanso
% Dall'ascesa delle casate al 1507 e oltre
% Rilevante per OpenHistoryMap come public history companion
%
% @ohm:layer "Genova 1507 — La Superba e le Guerre d'Italia"
% @ohm:timespan "1100/1512"
% @ohm:language "it"
% =============================================================================

\chapter{La Lunga Storia di Genova}
\label{ch:storia-espansa}

\begin{quote}
\itshape
Questo capitolo non è solo sfondo per il gioco. È la materia prima di un
evento di \textbf{public history}: venti giocatori che incarnano le grandi
famiglie genovesi del 1507 stanno anche esplorando, attraverso la fiction,
una delle crisi più acute della storia italiana del Rinascimento.
Le scelte che faranno questa sera rispecchiano quelle reali — e il
risultato reale è documentato.
\end{quote}

% =============================================================================
\section{Le Origini: dal Comune alla Repubblica (1099–1339)}
\label{sec:origini}
% =============================================================================

% @ohm:Event { "id": "genova:fondazione-comune", "date": "1099",
%   "label": "Fondazione del Comune di Genova",
%   "type": "ohm:PoliticalEvent", "place": "genova:CittaDiGenova" }

Genova diventa \textit{comune} nel \textbf{1099}, durante la Prima Crociata:
i mercanti genovesi finanziano le flotte crociate in cambio di quartieri
commerciali nelle città levantine. Questo è il modello genovese: non
conquista territoriale, ma \textbf{dominio dei nodi commerciali}.
Alessandria, Costantinopoli, Caffa in Crimea, Ceuta — ovunque ci sia uno
scalo strategico, c'è un \textit{fondaco} genovese.

\begin{multicols}{2}

\subsection{Guelfi e Ghibellini}

% @ohm:Actor { "id": "genova:fazione-guelfa", "label": "Fazione Guelfa",
%   "type": "ohm:PoliticalFaction", "timespan": "1200/1528" }
% @ohm:Actor { "id": "genova:fazione-ghibellina", "label": "Fazione Ghibellina",
%   "type": "ohm:PoliticalFaction", "timespan": "1200/1528" }

La divisione tra \idxbold{Guelfi} e \idxbold{Ghibellini} a Genova non è
meramente ideologica — è strutturale. I \textbf{Fieschi} guidano i Guelfi
(fedeli al Papa), i \textbf{Doria} e gli \textbf{Spinola} guidano i
Ghibellini (fedeli all'Imperatore). I \textbf{Grimaldi}, guelfi ma
pragmatici, oscillano secondo convenienza.

Questa divisione genera un ciclo di colpi di stato ogni decennio: quando
una fazione prevale, l'altra viene esiliata. Gli esiliati tornano.
Il ciclo ricomincia. Genova accumula ricchezza malgrado — o forse
\textit{grazie a} — questa instabilità permanente.

\columnbreak

\subsection{Il Banco di San Giorgio}

% @ohm:Actor { "id": "genova:banco-san-giorgio",
%   "label": "Banco di San Giorgio",
%   "type": "ohm:Institution",
%   "timespan": "1407/1805",
%   "place": "genova:PalazzoDiSanGiorgio" }
% @ohm:Place { "id": "genova:PalazzoDiSanGiorgio",
%   "label": "Palazzo di San Giorgio",
%   "geo:lat": 44.4088, "geo:long": 8.9269,
%   "ohm:accuracy": "high" }

Fondato nel \textbf{1407}, il \idxbold{Banco di San Giorgio} è forse la
prima banca pubblica moderna d'Europa. Gestisce il debito pubblico
genovese, controlla le colonie oltremare, e — crucialmente — è uno
\textbf{stato nello stato}: ha proprie leggi, proprie prigioni, propri
agenti diplomatici. Durante l'occupazione francese del 1507 rimane
formalmente neutro. Chiunque controlli Genova deve fare i conti con il
Banco — non il contrario.

\end{multicols}

% =============================================================================
\section{Il Dominio Visconteo e l'Età delle Crisi (1339–1453)}
\label{sec:crisi}
% =============================================================================

% @ohm:Event { "id": "genova:primo-doge", "date": "1339",
%   "label": "Simone Boccanegra primo Doge di Genova",
%   "type": "ohm:PoliticalEvent" }

Nel \textbf{1339} Genova istituisce il \textbf{Dogato} — un tentativo di
uscire dal ciclo fazioso eleggendo un doge a vita. Il primo è Simone
Boccanegra (la storia che Verdi renderà famosa quattrocento anni dopo).
Non funziona: i dogi vengono assassinati, esiliati, o deposti con
regolarità cronometrica.

\begin{multicols}{2}

Tra il 1353 e il 1421 Genova si mette volontariamente sotto la protezione
dei \textbf{Visconti di Milano} tre volte separate — ogni volta per
sfuggire a una crisi interna che le fazioni locali non sanno risolvere.
Questo schema — \textit{cedere la sovranità a una potenza esterna pur di
avere ordine} — è esattamente ciò che ha prodotto l'occupazione francese
del 1499.

\columnbreak

\begin{rulebox}[SharedBrown]{Nota per i Giocatori}
Il vostro personaggio del 1507 porta in sé questa memoria istituzionale.
Genova ha già \textit{scelto} di cedere la sovranità in passato.
La domanda non è \textit{se} è possibile — è \textit{a quali condizioni}
e \textit{con quali garanzie}. Questo spiega perché alcune famiglie
negoziano con i francesi anche mentre sostengono la rivolta.
\end{rulebox}

\end{multicols}

% =============================================================================
\section{L'Apogeo e la Caduta: 1453–1499}
\label{sec:apogeo}
% =============================================================================

% @ohm:Event { "id": "genova:caduta-caffa", "date": "1475",
%   "label": "Caduta di Caffa agli Ottomani",
%   "type": "ohm:MilitaryEvent",
%   "place": "ohm:Caffa" }
% @ohm:Event { "id": "genova:trattato-senlis", "date": "1493",
%   "label": "Trattato di Senlis — Francia rinuncia alla Borgogna",
%   "type": "ohm:DiplomaticEvent" }

La caduta di \textbf{Costantinopoli} (1453) e poi di \textbf{Caffa} (1475)
agli Ottomani decapita il sistema commerciale genovese nel Mediterraneo
orientale. Genova perde in un decennio ciò che aveva costruito in tre
secoli. Il vuoto economico è drammatico — e apre la strada alla
penetrazione dei creditori stranieri.

\begin{multicols}{2}

\subsection{Le Guerre d'Italia cominciano}

% @ohm:Event { "id": "genova:invasione-carlo-viii", "date": "1494",
%   "label": "Carlo VIII invade l'Italia",
%   "type": "ohm:MilitaryEvent" }

Nel \textbf{1494} Carlo~VIII di Francia attraversa le Alpi con 25.000
uomini — il più grande esercito che l'Italia abbia visto da secoli.
Conquista Napoli in sei mesi. L'equilibrio della penisola, costruito
dalla Pace di Lodi (1454) e mantenuto da Lorenzo de' Medici, crolla
in pochi mesi. Le \idxbold{Guerre d'Italia} (1494–1559) trasformeranno
la penisola in campo di battaglia tra Francia, Spagna e Impero per
sessant'anni.

\columnbreak

\subsection{Genova come pedina}

Genova è geograficamente cruciale: chiunque voglia controllare il nord
Italia deve controllare il suo porto. I milanesi Sforza — alleati
storici di alcune famiglie genovesi — vengono cacciati da Luigi~XII
nel \textbf{1499}. Con loro cade l'ultimo baluardo contro l'occupazione
diretta.

% @ohm:Event { "id": "genova:occupazione-francese", "date": "1499",
%   "label": "Luigi XII occupa il Ducato di Milano e Genova",
%   "type": "ohm:MilitaryEvent",
%   "place": "genova:CittaDiGenova" }

\end{multicols}

% =============================================================================
\section{L'Occupazione Francese (1499–1507)}
\label{sec:occupazione}
% =============================================================================

% @ohm:Actor { "id": "genova:luigi-xii",
%   "label": "Luigi XII re di Francia",
%   "type": "ohm:Person", "timespan": "1462/1515" }
% @ohm:Actor { "id": "genova:trivulzio",
%   "label": "Gian Giacomo Trivulzio",
%   "type": "ohm:Person", "timespan": "1441/1518",
%   "ohm:role": "Governatore di Genova 1499-1500, 1507" }

\begin{multicols}{2}

Sotto Luigi~XII, Genova non viene annessa alla Francia — viene trattata
come un \textbf{feudo personale del re}. La distinzione è importante:
significa che le sue istituzioni formali (il Doge, il Banco di San Giorgio,
i consigli delle famiglie) sopravvivono, ma senza potere reale.
Il governatore francese — prima \textbf{Trivulzio}, poi altri — tiene
la guarnigione nel castello di \textbf{Castelletto}, che domina la città
dall'alto.

\subsubsection{Il peso fiscale}

La tassazione francese è sistematica e pesante. Luigi~XII ha bisogno di
finanziare le sue campagne in Italia, e Genova è una delle sue fonti
principali. Le dogane vengono ri-tariffate, i monopoli commerciali
redistribuiti a favori di mercanti francesi, le esenzioni storiche
delle famiglie genovesi cancellate con decreti.

\columnbreak

\subsubsection{La guarnigione}

% @ohm:Place { "id": "genova:Castelletto",
%   "label": "Castello di Castelletto",
%   "type": "ohm:MilitaryStructure",
%   "geo:lat": 44.4120, "geo:long": 8.9275,
%   "ohm:accuracy": "approximate",
%   "ohm:note": "Demolito nel 1536. Posizione approssimativa." }

Il \idxbold{Castelletto} è il simbolo fisico dell'occupazione: una
fortezza medievale sul colle che sovrasta il porto, rinforzata dai
francesi e presidiata da una guarnigione permanente. È praticamente
inespugnabile dall'interno della città. Per questo il piano della
rivolta del 1507 non prevede di conquistarlo — ma di \textit{isolarlo}.

\begin{rulebox}[SharedBrown]{Nota geografica — OHM}
Il Castelletto fu demolito definitivamente nel 1536.
La sua posizione è ricostruibile da fonti cartografiche del XVI sec.
Layer OHM suggerito: \texttt{genova1507:military\_structures}.
\end{rulebox}

\end{multicols}

% =============================================================================
\section{La Rivolta del 1507: Cronologia Precisa}
\label{sec:rivolta-cronologia}
% =============================================================================

% @ohm:Event { "id": "genova:rivolta-1507",
%   "label": "Rivolta di Genova contro la Francia",
%   "type": "ohm:MilitaryEvent",
%   "timespan": "1507-01/1507-04",
%   "place": "genova:CittaDiGenova" }

Questa è la cronologia reale degli eventi che il vostro gioco simula
e comprime in una singola notte. È utile conoscerla — sia per i
giocatori che vogliono ancorare la fiction alla storia, sia per
capire \textit{perché} le scelte dei personaggi avevano un peso reale.

\begin{center}
\rowcolors{2}{LightBG}{white}
\begin{tabularx}{\textwidth}{L{3.2cm} X}
  \toprule
  \textbf{Data} & \textbf{Evento} \\
  \midrule
  \textbf{Estate 1506} &
    % @ohm:Event { "id":"genova:tensioni-1506","date":"1506-07","label":"Tensioni fiscali crescenti a Genova","type":"ohm:SocialEvent" }
    Le tensioni fiscali raggiungono il punto critico. I mercanti del porto
    rifiutano di pagare le nuove imposte sulle merci. Piccoli scontri tra
    popolani e guardie francesi. \\

  \textbf{Gennaio 1507} &
    % @ohm:Event { "id":"genova:primo-insurrezione","date":"1507-01","label":"Prima insurrezione popolare nei caruggi","type":"ohm:MilitaryEvent" }
    Prima insurrezione popolare nei \textit{caruggi}. I francesi rispondono
    con arresti e rappresaglie. Le famiglie nobili osservano, non intervengono
    ancora. \\

  \textbf{Febbraio 1507} &
    Le famiglie genovesi iniziano a coordinare — non pubblicamente, ma
    attraverso intermediari. Il governatore Trivulzio richiede rinforzi
    a Luigi~XII. \\

  \textbf{Aprile 1507} &
    % @ohm:Event { "id":"genova:rivolta-principale","date":"1507-04",
    %   "label":"Rivolta principale — le famiglie si schierano apertamente",
    %   "type":"ohm:MilitaryEvent" }
    La rivolta diventa aperta. Le famiglie nobiliari si schierano.
    Le campane di San Lorenzo suonano. Trivulzio si ritira nel Castelletto.
    Genova è de facto libera per alcune settimane. \\

  \textbf{Maggio 1507} &
    % @ohm:Event { "id":"genova:reconquista-francese","date":"1507-05",
    %   "label":"Luigi XII riconquista Genova","type":"ohm:MilitaryEvent" }
    Luigi~XII scende personalmente in Italia con un esercito. Genova viene
    riconquistata. Le rappresaglie sono severe ma selettive — il re ha
    bisogno dei mercanti genovesi per finanziare le sue campagne. \\

  \textbf{Giugno 1507} &
    % @ohm:Event { "id":"genova:amnistia","date":"1507-06",
    %   "label":"Amnistia parziale e nuovi accordi di governo","type":"ohm:DiplomaticEvent" }
    Amnistia parziale. Nuovi accordi di governo: Genova ottiene alcune
    concessioni fiscali in cambio della sottomissione formale.
    Le famiglie che hanno trattato — come i Grimaldi — ottengono
    condizioni migliori. \\
  \bottomrule
\end{tabularx}
\end{center}

% =============================================================================
\section{Dopo il 1507: Cosa Succede Davvero}
\label{sec:dopo}
% =============================================================================

\begin{multicols}{2}

Questa sezione rompe la quarta parete della fiction. I giocatori che
hanno appena vissuto la notte del 1507 meritano di sapere cosa è
accaduto davvero ai protagonisti storici.

\subsubsection{Andrea Doria — il personaggio che mancava}

% @ohm:Actor { "id": "genova:andrea-doria",
%   "label": "Andrea Doria",
%   "type": "ohm:Person", "timespan": "1466/1560",
%   "ohm:significance": "Ammiraglio, poi principe di Melfi" }

Il vero Andrea Doria nel 1507 ha quarantun anni ed è già un comandante
navale affermato. Combatte per i Fieschi, poi per il Papa, poi per
la Francia, poi per la Spagna. Nel \textbf{1528} cambia definitivamente
campo: porta Genova dalla parte di Carlo~V d'Asburgo, ottiene il titolo
di \textit{principe di Melfi}, e di fatto governa la Repubblica fino
alla morte nel \textbf{1560}. È il personaggio più importante della
storia genovese del XVI secolo — ed è intenzionalmente assente dalla
vostra sessione perché il suo destino era già scritto.

\columnbreak

\subsubsection{I Fieschi — la ribellione mancata}

% @ohm:Event { "id": "genova:congiura-fieschi", "date": "1547",
%   "label": "Congiura dei Fieschi contro i Doria",
%   "type": "ohm:MilitaryEvent" }

Nel \textbf{1547} Gian Luigi Fieschi — nipote del Niccolò del vostro
gioco — tenta un colpo di stato contro Andrea Doria. La congiura
fallisce per caso: Fieschi annegò cadendo in mare da una passerella
durante la notte del golpe. L'episodio ispira il \textit{Fiesco} di
Schiller (1783). I Fieschi perdono quasi tutto dopo il 1547.

\subsubsection{Gli Spinola — la continuità del potere}

Gli Spinola sopravvivono a tutto. Nel XVII e XVIII secolo sono tra
le famiglie più potenti d'Europa, con rami in Spagna, Fiandre, e
nell'Impero. Il loro modello — capitale finanziario, rete diplomatica,
adattamento al vincitore di turno — si rivela il più duraturo.

\subsubsection{I Grimaldi — Monaco oggi}

% @ohm:Actor { "id": "genova:grimaldi-monaco",
%   "label": "Casa Grimaldi di Monaco",
%   "type": "ohm:Dynasty", "timespan": "1297/" }

La casa Grimaldi governa il Principato di Monaco ancora oggi. È la
conseguenza diretta della politica di doppio gioco che il vostro
tavolo Grimaldi ha simulato questa sera. La loro strategia — mantenere
l'autonomia cedendo il minimo indispensabile alla potenza dominante —
funzionò per cinquecento anni.

\end{multicols}

% =============================================================================
\section{Genova come Spazio Geografico: I Sestieri nel 1507}
\label{sec:geografia}
% =============================================================================

% @ohm:Place { "id": "genova:CittaDiGenova",
%   "label": "Genova", "type": "ohm:City",
%   "geo:lat": 44.4056, "geo:long": 8.9463,
%   "ohm:accuracy": "high",
%   "ohm:timespan": "ancient/" }

Genova nel 1507 è organizzata in \textbf{sei sestieri}, ognuno con la
propria identità sociale, le proprie chiese di riferimento, e le proprie
famiglie dominanti. La topografia condiziona la tattica della rivolta:
chi controlla quale sestiere determina chi controlla la città.

\begin{center}
\rowcolors{2}{LightBG}{white}
\begin{tabularx}{\textwidth}{L{2.5cm} L{3cm} L{2.5cm} X}
  \toprule
  \textbf{Sestiere} & \textbf{Posizione} & \textbf{Famiglia dom.} & \textbf{Carattere} \\
  \midrule

  % @ohm:Place { "id": "genova:sestiere-pre",
  %   "label": "Sestiere di Prè", "type": "ohm:District",
  %   "geo:lat": 44.4080, "geo:long": 8.9170,
  %   "ohm:accuracy": "approximate",
  %   "ohm:partOf": "genova:CittaDiGenova" }
  \textbf{Prè} & Ovest, porto vecchio & Doria & Marinai, arsenali, pescatori. Cuore della resistenza popolare. \\

  % @ohm:Place { "id": "genova:sestiere-portoria",
  %   "label": "Sestiere di Portoria", "type": "ohm:District",
  %   "geo:lat": 44.4065, "geo:long": 8.9420,
  %   "ohm:accuracy": "approximate" }
  \textbf{Portoria} & Centro-est & Fieschi & Artigiani, botteghe, la base popolare dei Fieschi. \\

  % @ohm:Place { "id": "genova:sestiere-san-giovanni",
  %   "label": "Sestiere di San Giovanni", "type": "ohm:District",
  %   "geo:lat": 44.4095, "geo:long": 8.9280,
  %   "ohm:accuracy": "approximate" }
  \textbf{San Giovanni} & Centro, porto & Spinola & Mercanti, banchieri, il cuore finanziario. \\

  % @ohm:Place { "id": "genova:sestiere-san-lorenzo",
  %   "label": "Sestiere di San Lorenzo", "type": "ohm:District",
  %   "geo:lat": 44.4071, "geo:long": 8.9320,
  %   "ohm:accuracy": "approximate" }
  \textbf{San Lorenzo} & Centro storico & Fieschi & La cattedrale, il palazzo ducale, il potere civile e religioso. \\

  % @ohm:Place { "id": "genova:sestiere-castello",
  %   "label": "Sestiere di Castello", "type": "ohm:District",
  %   "geo:lat": 44.4060, "geo:long": 8.9380,
  %   "ohm:accuracy": "approximate" }
  \textbf{Castello} & Est, collina & Doria/Spinola & Nobiltà mista, palazzi, posizione difensiva elevata. \\

  % @ohm:Place { "id": "genova:sestiere-soziglia",
  %   "label": "Sestiere di Soziglia", "type": "ohm:District",
  %   "geo:lat": 44.4082, "geo:long": 8.9350,
  %   "ohm:accuracy": "approximate" }
  \textbf{Soziglia} & Centro-nord & Grimaldi & Commercio internazionale, fondachi stranieri, zona di transito. \\
  \bottomrule
\end{tabularx}
\end{center}

\begin{multicols}{2}

\subsection{I Luoghi Chiave della Rivolta}

% @ohm:Place { "id": "genova:SanLorenzo",
%   "label": "Cattedrale di San Lorenzo",
%   "type": "ohm:ReligiousBuilding",
%   "geo:lat": 44.4071, "geo:long": 8.9316,
%   "ohm:accuracy": "high",
%   "ohm:significance": "Centro spirituale e politico della rivolta" }

\textbf{Cattedrale di San Lorenzo} — Le campane di San Lorenzo sono
il segnale della rivolta. La cattedrale è anche il luogo dove si tengono
le assemblee civiche informali. Chi controlla San Lorenzo controlla il
\textit{simbolo} della città.

% @ohm:Place { "id": "genova:PalazzoDucale",
%   "label": "Palazzo Ducale",
%   "type": "ohm:GovernmentBuilding",
%   "geo:lat": 44.4068, "geo:long": 8.9337,
%   "ohm:accuracy": "high" }

\textbf{Palazzo Ducale} — Sede del governo formale. Nel 1507 è
occupato dall'amministrazione francese. Chi lo controlla alla fine
della rivolta tiene in mano il potere esecutivo.

\columnbreak

% @ohm:Place { "id": "genova:Porto",
%   "label": "Porto di Genova",
%   "type": "ohm:Harbour",
%   "geo:lat": 44.4063, "geo:long": 8.9237,
%   "ohm:accuracy": "high" }

\textbf{Porto} — Il cuore economico. Le galee dei Doria vi sono
ancorate. Chi controlla il porto controlla i rifornimenti e la fuga
— sia dei francesi che dei rivoluzionari sconfitti.

% @ohm:Place { "id": "genova:Molo",
%   "label": "Il Molo (molo principale)",
%   "type": "ohm:Infrastructure",
%   "geo:lat": 44.4046, "geo:long": 8.9268,
%   "ohm:accuracy": "approximate" }

\textbf{Il Molo} — Il molo principale, punto di controllo per
tutte le merci in entrata e uscita. Tassato dai francesi, presidiato
da una guarnigione mista.

\end{multicols}

% =============================================================================
\section{Genova come Evento di Public History}
\label{sec:public-history}
% =============================================================================

\begin{rulebox}[GoldRPG]{Per Organizzatori e Facilitatori}
Questo gioco è progettato per funzionare come \textbf{evento di public
history} — un formato in cui la partecipazione ludica diventa un veicolo
per la comprensione storica. Alcune indicazioni pratiche:

\medskip
\textbf{Prima della sessione (30 minuti):}
Illustrare brevemente il contesto storico usando la cronologia in
\S\ref{sec:rivolta-cronologia}. Non è necessario che i giocatori
conoscano ogni dettaglio — l'obiettivo è dare una \textit{sensazione
di autenticità} che rende le scelte significative.

\medskip
\textbf{Durante la sessione:}
I GM possono introdurre dettagli storici come informazioni in-game:
un PNG che menziona le perdite commerciali dopo Caffa,
una lettera che cita le tasse di Luigi~XII. La fiction e la storia
si rinforzano a vicenda.

\medskip
\textbf{Dopo la sessione (45 minuti):}
Debrief strutturato in due parti:
\begin{enumerate}
  \item \textit{Cosa è successo nel gioco?} — ogni tavolo racconta il proprio epilogo.
  \item \textit{Cosa è successo davvero?} — confronto con la cronologia reale
        (\S\ref{sec:dopo}). I giocatori scoprono cosa è accaduto ai
        ``loro'' personaggi storici nei decenni successivi.
\end{enumerate}

\medskip
\end{rulebox}
